\chapter{1983 Nobel Prize in Economics}
\textbf{Laureates: }Gérard Debreu (France/USA)

\section*{Reason for Award}
For his having incorporated new analytical methods into economic theory and for his rigorous reformulation of the theory of general equilibrium

\section{What They Did}
Gérard Debreu provided the mathematical proof of the existence of competitive equilibrium using convex set theory, transforming general equilibrium into a mathematically rigorous discipline. His seminal work Theory of Value established general equilibrium theory on proper mathematical foundations, utilizing topological fixed-point theorems to demonstrate existence of equilibrium under well-defined assumptions. He showed that under convex preferences and continuous functions with well-behaved production sets, competitive equilibrium exists. Debreu contributed to utility theory, proving existence of continuous utility functions representing ordinal preferences. He provided rigorous treatment of arbitrage-free pricing in complete markets, establishing the foundation for modern asset pricing theory. His work brought mathematical precision to economics, allowing precise statements about existence, uniqueness, and stability of equilibria.

\section{Impact on Economic Thinking}
Debreu's work transformed general equilibrium from heuristic reasoning into a rigorous mathematical discipline, enabling precise statements about existence, uniqueness, and stability of equilibria. His formulation of arbitrage-free pricing provided the foundation for modern asset pricing theory. By establishing general equilibrium under convexity conditions, he clarified when markets achieve Pareto optimality, giving precise conditions under which competitive allocations are socially optimal. His mathematical approach influenced generations of theorists who extend his framework to incomplete markets, dynamic settings, and international trade. Debreu demonstrated that rigorous mathematics could address fundamental economic questions about market efficiency and welfare, elevating the standards of economic proof and argument.

\section{Modern Perspectives Today}
General equilibrium theory remains a cornerstone of theoretical economics, with Debreu's formulation providing essential mathematical tools. Existence, uniqueness, and stability analyses build upon his foundational work. Asset pricing theory relies on his arbitrage-free characterization of security values. Extensions to incomplete markets, dynamic general equilibrium, search frictions, and heterogeneous agents all depend on the mathematical clarity Debreu established. He brought precision to economics, allowing researchers to explore complex questions about market failures, public goods, and externalities with rigorous methodology. His work underpins much of contemporary microeconomic theory and continues to guide research directions in general equilibrium analysis, mathematical economics, and financial economics.
